% Bagian: propositional-logic
% Bab: propositional-logic
% Subbagian: soundness

\documentclass[../../../include/open-logic-section]{subfiles}

\begin{document}

\olfileid[id]{pl}{prp}{snd}

\olsection{Kesahihan Logika Proposisi}

\begin{thm}[Kesahihan]
\ollabel{thm:soundness}
Jika $\Gamma\Proves!A$, maka
$\Gamma \Entails !A$.
\end{thm}

\begin{proof}
Melalui induksi pada derivasi. Jika $!A$ merupakan suatu aksioma, maka
$\Entails !A$ dan, terlebih lagi, $\Gamma \Entails!A$. Demikian pula, jika
$!A \in \Gamma$, maka $\Gamma \Entails!A$. Untuk langkah induktif, misalkan
$!A$ diperoleh melalui \emph{modus ponens} dari $!C$ dan $!C\lif!A$; misalkan
pula, sebagai hipotesis induksi, bahwa teorema tersebut berlaku bagi
$!C\lif !A$ dan $!C$. Butir ketiga dalam
\olref[sem]{prop:semanticalfacts} memberikan hasil yang dikehendaki.
\end{proof}

\begin{cor}
Jika $\Gamma$ dapat dipenuhi, maka $\Gamma$ konsisten. Jadi, logika proposisi
konsisten.
\end{cor}

\end{document}
